\documentclass[11pt]{article}
\usepackage[T1]{fontenc}
\usepackage{lmodern,amsmath,amssymb,amsthm,mathtools}
\usepackage[margin=1in]{geometry}
\usepackage[colorlinks=true,linkcolor=blue,citecolor=blue,urlcolor=blue]{hyperref}
\usepackage{microtype}
\newtheorem{theorem}{Theorem}[section]
\newtheorem{lemma}[theorem]{Lemma}
\newtheorem{proposition}[theorem]{Proposition}
\newtheorem{corollary}[theorem]{Corollary}
\theoremstyle{remark}\newtheorem{remark}[theorem]{Remark}
\DeclareMathOperator{\per}{per}
\DeclareMathOperator{\Crit}{Crit}
\DeclareMathOperator{\codim}{codim}
\DeclareMathOperator{\Gr}{Gr}
\DeclareMathOperator{\dc}{dc}
\DeclareMathOperator{\rank}{rank}
\DeclareMathOperator{\adj}{adj}
\newcommand{\bdc}{\overline{\dc}}
\newcommand{\C}{\mathbb C}
\newcommand{\PP}{\mathbb P}
\title{Smooth slices and superquadratic border determinantal\\
complexity of the permanent}
\author{Ying Xie\\Kennesaw State University\\
\texttt{yxie2@kennesaw.edu}}
\date{}
\ifdefined\pdfinfoomitdate\pdfinfoomitdate=1\fi
\ifdefined\pdfsuppressptexinfo\pdfsuppressptexinfo=15\fi
\ifdefined\pdftrailerid\pdftrailerid{}\fi
\hypersetup{pdftitle={Smooth slices and superquadratic border determinantal complexity of the permanent},
pdfauthor={Ying Xie},pdfsubject={Algebraic complexity and polar intersections}}
\hypersetup{pdfcreator={},pdfproducer={},pdfcreationdate={},pdfmoddate={}}
\ifdefined\XeTeXversion
\AtBeginDocument{\special{pdf:docinfo << /CreationDate () /ModDate () /Creator () /Producer () >>}
\special{pdf:trailerid [<00000000000000000000000000000000> <00000000000000000000000000000000>]}}
\fi
\begin{document}
\maketitle
\begin{abstract}
We prove that the border determinantal complexity of the $n\times n$
permanent over $\mathbb C$ is $\Omega(n^2\log n)$. In particular, it is
at least $n^2(\log_2 n-5)/(64e)$ for $n\ge64$, and at least
$3n^2(\log_2 n-8)/(256eH_2(1/8))$ for $n\ge512$, where $H_2$ is
binary entropy. The proof combines the matching-polynomial restrictions
of OpenAI with Sheshadri's determinantal kernel incidence, using a
general linear slice to obtain a smooth homogeneous target. A direct
perturbation argument transfers its polar count to coefficient limits
of arbitrary affine-linear determinants. We refine the shared matching
coefficients by removing their constant block on the critical locus;
the remaining coefficient compression also appears in work of Guo.
The argument gives corresponding bounds for arbitrary acyclic algebraic
branching programs and determines the order of border determinantal
complexity of elementary symmetric polynomials. An algebraic transfer
retains multiplicities and extends the permanent bounds to algebraically
closed fields of characteristic $p>n$. Separate circuit-state bounds
allow $O(n^2)$ non-skew products and nonscalar divisions. These restricted
circuit consequences do not separate VP from VNP.
\end{abstract}

\section{Introduction and main results}

For a polynomial $f$, its border determinantal complexity $\bdc(f)$
is the least integer $m$ for which $f$ belongs to the Zariski closure of
\[
\mathcal D_m=\left\{\det\left(A_0+\sum_i x_i A_i\right):
 A_i\in\C^{m\times m}\right\}
\]
in the space of polynomials of degree at most $m$. This agrees with
coefficientwise approximation over $\C$: $\mathcal D_m$ is constructible,
and a complex constructible set has the same Euclidean and Zariski
closures. Exact determinantal complexity is denoted by $\dc$.

Landsberg, Manivel, and Ressayre~\cite{LMR} proved the quadratic bound
$\bdc(\per_n)\ge n^2/2$. We obtain an unbounded logarithmic
improvement in the asymptotic lower bound.

\begin{theorem}\label{thm:main}
Over $\mathbb C$, for every integer $n\ge64$,
\begin{equation}\label{eq:main-simple}
 \bdc(\per_n)\ge\frac{n^2(\log_2 n-5)}{64e}.
\end{equation}
Writing $\gamma=H_2(1/8)$, for every $n\ge512$ one also has
\begin{equation}\label{eq:main-compressed}
 \bdc(\per_n)\ge\frac{3n^2(\log_2 n-8)}{256e\gamma}.
\end{equation}
In particular,
\[
 \liminf_{n\to\infty}\frac{\bdc(\per_n)}{n^2\log_2 n}
 \ge\frac{3}{256eH_2(1/8)}.
\]
Here $H_2(x)=-x\log_2 x-(1-x)\log_2(1-x)$.
The same bounds hold for exact determinantal complexity.
\end{theorem}

These explicit inequalities need not exceed $n^2/2$ throughout their
stated finite ranges; either may always be combined with the quadratic
bound. The first has a shorter parameter calculation, while the second
has the larger asymptotic constant.

The proof combines the matching-polynomial restriction and critical-locus
bound of~\cite[Chapter 5, Proposition 5.1 and Lemma 6.1]{OAI} with the
determinantal polar-incidence mechanism of~\cite[Theorem 1.3]{Sheshadri}.
We include derivations of the needed matching identities and a direct
smooth-target transfer proof, so the main argument does not depend on
other claims in those works. Guo~\cite[Proposition 3.1]{Guo} obtains the
same disjointness-matrix compression of the shared matching coefficients.
We further use the vanishing of their constant block on the critical
locus. Neither the matching selector nor the kernel incidence is claimed
as a new construction here. The main result is their composition through
a fixed smooth slice, with explicit finite-size estimates.

The elementary-symmetric consequence is in Section~\ref{sec:elementary}.
Appendix~\ref{app:algebraic} gives the multiplicity-preserving transfer
and its large-characteristic consequence. Appendices~\ref{app:nonskew}
and~\ref{app:rationalstate} concern circuits with a stated budget on
non-skew products and nonscalar divisions. The latter statements keep
that budget and are distinct from unrestricted circuit lower bounds,
such as those in~\cite{OAI,Guo}.

All critical loci in this note are affine zero sets of all first
partial derivatives. A smooth projective hypersurface means a smooth
hypersurface scheme, not merely a smooth reduced support. Every degree
under consideration is at least two.

\section{From critical codimension to a smooth slice}

\begin{lemma}\label{lem:slice}
Let $P$ be a nonzero homogeneous degree-$d$ form on $\C^a$, where
$d\ge2$. Suppose $\codim\Crit(P)\ge k\ge3$. Then a general $k$-dimensional
linear subspace $W\subseteq\C^a$ has the property that
$V(P|_W)\subset\PP(W)$ is a smooth degree-$d$ hypersurface.
\end{lemma}

\begin{proof}
If $k=a$, the critical locus is $\{0\}$ and there is nothing to prove.
Suppose $k<a$. The Grassmannian $\Gr(k,a)$ has dimension $k(a-k)$.
For a smooth point $[x]\in V(P)\subset\PP^{a-1}$, the restriction can
be singular at $[x]$ only if
\[
 [x]\in\PP(W)\subseteq T_{[x]}V(P).
\]
For this fixed point, the choices of $W$ form a Grassmannian of
dimension $(k-1)(a-k-1)$. Since the smooth locus of the hypersurface
has dimension $a-2$, the incidence of these pairs has dimension
\[
 (a-2)+(k-1)(a-k-1)=k(a-k)-1.
\]
The projective singular locus has dimension at most $a-k-1$, by
homogeneity and Euler's identity. For a fixed singular point, the
subspaces containing it have dimension $(k-1)(a-k)$. Thus the other
bad incidence also has dimension at most $k(a-k)-1$. Their images
cannot fill the Grassmannian. A general $W$ avoids both and has nonzero
restriction, proving the claim.
\end{proof}

\begin{remark}
This argument is necessary: restricting the original gradient and
projecting its output is not the same operation as taking the gradient
of a restriction. The tangent incidence excludes singularities created
by the latter operation. The subspace $W$ will be fixed before any
border limit is taken.
\end{remark}

\section{A border bound for smooth homogeneous targets}

We give the precise smooth-target consequence of the kernel-incidence
method~\cite{Sheshadri} needed here. The proof uses persistence of simple
affine polar points, avoiding a general assertion about specialization
of conormal cycles.

\begin{proposition}\label{prop:polar}
Let $f$ be a homogeneous degree-$d$ form in $k\ge3$ variables, with
$V(f)\subset\PP^{k-1}$ smooth. If $\bdc(f)\le m$, then
\begin{align}
 d(d-1)^{k-2}&\le B(m,k),\label{eq:polar}\\
 B(m,k)&=\sum_{a=1}^{k-1}
 \binom ma\binom{m-1}{k-1-a}\binom{k-2}{a-1}.
 \label{eq:B}
\end{align}
Consequently,
\begin{equation}\label{eq:bridge}
 \bdc(f)\ge\frac{(k-1)(d-1)}{4e}.
\end{equation}
\end{proposition}

\begin{proof}
\emph{Simple points on the target.}
Euler's identity and smoothness imply that the first partials of $f$
have no common projective zero. They define a morphism on
$\PP^{k-1}$ whose pullback of $\mathcal O(1)$ is
$\mathcal O(d-1)$. This morphism is finite: a positive-dimensional
fiber would contain a curve on which an ample line bundle had degree
zero. Its restriction to $V(f)$ is finite as well.

Choose $k-2$ general constant-coefficient directional derivatives
$D_1,\ldots,D_{k-2}$. In characteristic zero a general linear section
of the image avoids the branch locus and pulls back transversely.
The intersection on $V(f)$ therefore consists of
\[
 \delta=d(d-1)^{k-2}
\]
distinct points. Equivalently, this follows by successive general
members of the basepoint-free polar system and Bezout. Choose a linear
form $\ell$ nonzero at all these points and rescale each representative
to satisfy $\ell=1$. The resulting points are simple zeros of the
square affine system
\begin{equation}\label{eq:affine}
 f=0,\qquad \ell-1=0,\qquad D_jf=0\quad(1\le j\le k-2).
\end{equation}

\emph{Persistence under approximation.}
Work in the coefficient space of polynomials of degree at most $m$.
The implicit-function theorem at each of the finitely many simple
zeros in~\eqref{eq:affine} gives a neighborhood of $f$ in which every
polynomial $g$ has a distinct simple zero near each one for the system
with $f$ replaced by $g$. Since $f\in\overline{\mathcal D_m}$, choose
$g=\det A$ in this neighborhood, where $A$ is an affine-linear
$m\times m$ matrix. At the retained zeros, $\nabla g\ne0$, so
$\rank A=m-1$. No convergence of the coefficients of $A$ is required.

\emph{A square kernel-incidence system.}
Write $\widehat A(z,x)=zA_0+\sum_i x_iA_i$. For each retained point,
choose its unique projective left and right kernel lines, represented
by $u$ and $v$. Choose a linear map
$\Lambda:\C^m\to\C^{m-1}$ injective on the image of $A$ at every
one of these finitely many points. Each is a nonempty open condition
on $\Lambda$, so a common choice exists. On
\[
 \PP^k_{[z:x]}\times\PP^{m-1}_{[u]}\times\PP^{m-1}_{[v]}
\]
consider the following system:
\begin{equation}\label{eq:incidence}
 u^{\mathsf T}\widehat A=0,\qquad
 \Lambda\widehat A v=0,\qquad
 u^{\mathsf T}(D_jA)v=0,\qquad
 \ell(x)-z=0.
\end{equation}
There are respectively $m$, $m-1$, $k-2$, and one equations. Their
multidegrees are $(1,1,0)$, $(1,0,1)$, $(0,1,1)$, and $(1,0,0)$.
The number of equations equals the ambient dimension.

\emph{Isolation is local, not a genericity assumption on $A$.}
Fix a retained point and work in the affine chart $z=1$. After
permuting rows and columns, write
\[
 A=\begin{pmatrix}B&c\\ r&s\end{pmatrix},\qquad \det B\ne0.
\]
In suitable kernel charts the full left equations solve
$u^{\mathsf T}=(-rB^{-1},1)$ and impose $s-rB^{-1}c=0$.
On this hypersurface, injectivity of $\Lambda$ on the matrix image
persists locally. The reduced right equations thus solve
$v=(-B^{-1}c,1)$. In particular, the kernel incidence is locally a
graph over $g=0$. On that graph,
\[
 \adj(A)=(\det B)vu^{\mathsf T},\qquad
 D_jg=(\det B)u^{\mathsf T}(D_jA)v.
\]
The factor $\det B$ is a unit. Thus the retained simple zeros of the
affine polar system lift to reduced isolated solutions
of~\eqref{eq:incidence}. Other components, including components at
infinity or of higher corank, do not affect this assertion.

\emph{Counting.}
Multihomogeneous Bezout bounds the number of isolated solutions by
\begin{equation}\label{eq:coefficient}
 [H^kU^{m-1}V^{m-1}]
 H(H+U)^m(H+V)^{m-1}(U+V)^{k-2}.
\end{equation}
This remains a bound when there are excess components: perturb the
sections of the globally generated line bundles to a proper
intersection; each reduced isolated solution persists. Expanding the
coefficient in~\eqref{eq:coefficient} gives~\eqref{eq:B}, and proves
\eqref{eq:polar}.

Finally, Vandermonde's identity gives
\[
 B(m,k)\le 2^{k-2}\binom{2m-1}{k-1}
 \le 2^{k-2}\left(\frac{e(2m-1)}{k-1}\right)^{k-1}.
\]
If the binomial coefficient vanishes,~\eqref{eq:polar} is impossible;
otherwise the displayed estimates imply
\[
 m\ge\frac12+\frac{k-1}{4e}
       \bigl(2d(d-1)^{k-2}\bigr)^{1/(k-1)}
 \ge\frac{(k-1)(d-1)}{4e}.
\]
\end{proof}

\begin{corollary}\label{cor:codim}
If $P$ is homogeneous of degree $d\ge2$ and
$\codim\Crit(P)\ge k\ge3$, then
\[
 \bdc(P)\ge\frac{(k-1)(d-1)}{4e}.
\]
The same lower bound holds for any polynomial having a nonzero scalar
multiple of $P$ as an affine restriction.
\end{corollary}
\begin{proof}
Apply Lemma~\ref{lem:slice} and Proposition~\ref{prop:polar}.
Affine substitution preserves size-$m$ determinantal representations
and commutes with coefficientwise limits. A nonzero scalar is absorbed
by scaling one row of the representing matrix.
\end{proof}

\section{The permanent specialization and its parameters}

We state the matching inputs and their coefficient refinement.
Appendix~\ref{app:matching} supplies all the required identities.
For a $t\times s$ matrix
of variables $X$, let
\[
 M_{t,s,d}(X)=\sum_{\substack{I\subseteq[t],\ J\subseteq[s]\\
                              |I|=|J|=d}}
                 \per(X_{I,J}).
\]
Proposition 5.1 of~\cite[Chapter 5]{OAI} states, for
$3\le d\le\min(t,s)$,
\begin{equation}\label{eq:critical}
 \dim\Crit(M_{t,s,d})\le 2^t-1+s(d-2).
\end{equation}
The refined estimate proved in Appendix~\ref{app:compression} is
\begin{equation}\label{eq:critical-compressed}
 \dim\Crit(M_{t,s,d})\le R_{\rm crit}(t,d)+s(d-2),\qquad
 R_{\rm crit}(t,d)=\sum_{j=1}^{d-1}\binom{t}{\min(j,d-j)}-t.
\end{equation}
For $b$ disjoint blocks of variables $X^{(h)}$ and nonzero scalars
$\lambda_h$, put
\[
P=\sum_{h=1}^b\lambda_h M_{t,s,d}(X^{(h)}).
\]
Its critical locus is the product of the
individual critical loci. Consequently,
\begin{equation}\label{eq:codimension}
 \codim\Crit(P)\ge b\bigl(s(t-d+2)-2^t+1\bigr).
\end{equation}

Lemma 6.1 of the same chapter supplies suitable nonzero $\lambda_h$
such that $P$ is a nonzero scalar multiple of an affine restriction
of $\per_{bt+s-d}$. More explicitly, if $r=bt$, partition the $r$
rows into blocks $R_1,\ldots,R_b$ of size $t$. For a primitive $d$th
root $\zeta$, take the constant columns of $U$ to be $t-d$ copies of
$\mathbf1_{R_h}$ for each $h$, and
\[
 \mathbf1_{R_1\cup\cdots\cup R_{h-1}}+2\zeta^j\mathbf1_{R_h}
 \quad(2\le h\le b,\ 0\le j<d).
\]
Then $U$ has $r-d$ columns, and the required affine matrix is
\[
 \begin{pmatrix}X&U\\ \mathbf1&0\end{pmatrix},
\]
where the lower block has $s-d$ rows. The common nonzero scalar is
$(s-d)!(t-d)!(t!)^{b-1}$. With
$\theta=(-1)^{d+1}2^d$, the weights may be written
\[
 \lambda_1=\theta^{b-1},\qquad
 \lambda_h=\theta^{b-h}(1+\theta)^{h-2}\quad(2\le h\le b).
\]
In particular no weight vanishes.

\begin{proposition}[Parameter bound]\label{prop:parameters}
Let $b\ge1$ and $3\le d\le\min(t,s)$ be integers, with
$n_0=bt+s-d\le n$. Set
\[
 K=b\bigl(s(t-d+2)-R_{\rm crit}(t,d)\bigr).
\]
If $K\ge3$, then
\[
 \bdc(\per_n)\ge\frac{(K-1)(d-1)}{4e}.
\]
One may instead use $K=b(s(t-d+2)-2^t+1)$.
\end{proposition}
\begin{proof}
The product description of the critical locus and
\eqref{eq:critical-compressed} give codimension at least $K$.
Use the selector above, Corollary~\ref{cor:codim}, and, when $n_0<n$,
pad the permanent specialization with an identity block.
\end{proof}

\subsection{Explicit parameter choices}

\begin{proof}[Proof of Theorem~\ref{thm:main}]
We first prove~\eqref{eq:main-simple} using the uncompressed estimate.
For $n\ge64$, set
\[
 t=\lfloor\log_2 n\rfloor,\qquad d=\lfloor t/2\rfloor,\qquad
 b=\left\lfloor\frac{n}{2t}\right\rfloor,\qquad
 r=bt,\qquad s=n-r+d.
\]
Then $t\ge6$, $b\ge1$, $3\le d\le\min(t,s)$, and $r+s-d=n$.
Since $n/2-t<r\le n/2$,
\[
 rs=r(n-r+d)\ge\frac{n^2}{4}-t^2.
\]
For $K=b(s(t-d+2)-2^t+1)$, using $2^t\le n$ gives
\begin{align*}
 K&=rs\frac{t-d+2}{t}-\frac rt(2^t-1)\\
  &\ge\left(\frac{n^2}{4}-t^2\right)
          \left(\frac12+\frac2t\right)-\frac{n^2}{2t}
   =\frac{n^2}{8}-\frac{t^2}{2}-2t.
\end{align*}
This lower bound is greater than three: at $t=6$ it is at least
$512-18-12$, and $2^{2t}/8-t^2/2-2t$ increases thereafter.
Also $K\le rs$. As $d-1\ge(t-3)/2$, we have
\[
 (K-1)(d-1)\ge
 \frac{n^2(t-3)}{16}
 -\frac{(t^2/2+2t+1)(t-3)}2.
\]
For $t\ge6$,
\[
 8(t^2/2+2t+1)(t-3)\le8t^3\le2^{2t}\le n^2.
\]
For the middle inequality, check $t=6$ and note that
$((t+1)/t)^3<4$. Consequently
\[
 (K-1)(d-1)\ge\frac{n^2(t-4)}{16}
 >\frac{n^2(\log_2 n-5)}{16}.
\]
Proposition~\ref{prop:parameters} proves~\eqref{eq:main-simple}.

For~\eqref{eq:main-compressed}, put
\[
 \gamma=H_2(1/8),\qquad c=\frac{8^8}{7^7}=2^{8\gamma}.
\]
Let $a$ be the largest integer with $c^a\le n$, and set
\[
 t=8a,\qquad d=2a,\qquad b=\left\lfloor\frac{n}{16a}\right\rfloor,
 \qquad r=bt,\qquad s=n-r+d.
\]
Since $20<c<21$ and $n\ge512$, we have $a\ge2$ and $b\ge1$;
all hypotheses of Proposition~\ref{prop:parameters} hold.
The profile in Appendix~\ref{app:compression} gives
\begin{equation}\label{eq:entropy-cost}
 R_{\rm crit}(8a,2a)
 =2\sum_{j=1}^{a-1}\binom{8a}{j}+\binom{8a}{a}-8a
 \le\frac43\binom{8a}{a}\le\frac43c^a\le\frac43n.
\end{equation}
Indeed, the ratios of consecutive terms below $\binom{8a}{a}$
are at most $a/(7a+1)<1/7$. For the next inequality, the single
$a$th term in the binomial expansion of $(1/8+7/8)^{8a}$ is at
most one, so $\binom{8a}{a}\le (8^8/7^7)^a$.

Write $K=b(s(t-d+2)-R_{\rm crit}(t,d))$. As before,
$rs\ge n^2/4-64a^2$; therefore
\begin{align*}
 K&\ge\left(\frac{n^2}{4}-64a^2\right)
              \left(\frac34+\frac{1}{4a}\right)-\frac{n^2}{12a}\\
  &=\frac{3n^2}{16}-\frac{n^2}{48a}-48a^2-16a.
\end{align*}
For $a\ge2$, the polynomial error satisfies
\[
 (48a^2+16a+1)(2a-1)\le96a^3,\qquad
 48\cdot96a^3\le20^{2a}<c^{2a}\le n^2.
\]
The second inequality holds at $a=2$ and is preserved on increasing
$a$, because $((a+1)/a)^3<400$. It follows that
\begin{align*}
 (K-1)(2a-1)
 &\ge n^2\left(\frac{3a}{8}-\frac{11}{48}
                         +\frac{1}{48a}\right)
      -(48a^2+16a+1)(2a-1)\\
 &\ge n^2\left(\frac{3a}{8}-\frac14\right).
\end{align*}
This estimate in particular ensures $K\ge3$; for example its
right side divided by $2a-1$ is greater than $n^2/8$.
Also $K\le rs$, directly from its definition.
Since $a>\log_2 n/(8\gamma)-1$ and $\gamma<3/5$,
\[
 (K-1)(d-1)
 \ge n^2\left(\frac{3\log_2 n}{64\gamma}-\frac58\right)
 \ge\frac{3n^2(\log_2 n-8)}{64\gamma}.
\]
Another application of Proposition~\ref{prop:parameters} proves
\eqref{eq:main-compressed}. Finally $\dc\ge\bdc$.
\end{proof}


\begin{corollary}[Unlayered algebraic branching programs]\label{cor:abp}
Let an algebraic branching program mean an arbitrary finite directed
acyclic graph with a designated source and sink and affine-linear
edge labels over $\C$. Its output is the sum of the weights of all
source-to-sink paths. If $\per_n$ is a coefficientwise limit of outputs
of such programs with at most $v$ vertices, then
\[
 v+1\ge \frac{n^2(\log_2 n-5)}{64e}\quad(n\ge64),
\]
and
\[
 v+1\ge \frac{3n^2(\log_2 n-8)}{256eH_2(1/8)}\quad(n\ge512).
\]
The same bound holds for exact computation. No layering, homogeneity
of intermediate computations, or variable-order restriction is assumed.
\end{corollary}

\begin{proof}
In a topological ordering, let $A$ be the weighted adjacency matrix
of the program; parallel-edge labels can be added. It is strictly
upper triangular, and its path polynomial is
\[
 e_s^{\mathsf T}(I-A)^{-1}e_t
   =\det\begin{pmatrix}I-A&e_t\\-e_s^{\mathsf T}&0\end{pmatrix}.
\]
This is an affine-linear determinant of size at most $v+1$. Pad with
identity blocks if necessary and take coefficientwise limits.
Theorem~\ref{thm:main} applies.
\end{proof}

The model distinction matters at polynomial scales. For example,
Chatterjee, Kumar, She, and Volk~\cite[Theorems 1.2 and 1.3]{CKSV}
obtain different bounds for layered and unlayered branching programs.
Corollary~\ref{cor:abp} is stated directly for arbitrary acyclic graphs.
It is an immediate consequence of the determinantal statement; it is
not a simulation of arbitrary arithmetic circuits by these graphs.

\section{An elementary-symmetric control family}\label{sec:elementary}

The same bridge has a consequence independent of the matching
specialization. It provides a useful check on both the degree dependence
and the distinction between circuits and determinantal representations.
Write $e_{N,d}$ for the elementary symmetric polynomial of degree $d$
in $N$ variables.

\begin{corollary}\label{cor:elementary}
Uniformly for $2\le d\le N$, over $\C$,
\[
 \bdc(e_{N,d})=\Theta\bigl(d(N-d+1)\bigr),\qquad
 \dc(e_{N,d})=\Theta\bigl(d(N-d+1)\bigr).
\]
For $d\le N-1$, a lower bound is
$(N-d+1)(d-1)/(4e)$. An exact upper bound is
$(d+1)(N-d+1)+1$.
\end{corollary}

\begin{proof}
The critical locus of $e_{N,d}$ consists of vectors with at most
$d-2$ nonzero coordinates; this is a known characteristic-zero
fact, recalled in~\cite{Orzel}. Here is a direct proof.
The identities
\[
 \sum_i\partial_i e_{N,d}=(N-d+1)e_{N,d-1},\qquad
 \partial_i e_{N,d}=e_{N,d-1}-x_i\partial_i e_{N,d-1}
\]
imply that, at a critical point, all derivatives of $e_{d-1}$ on
its nonzero support vanish. If the support has size $p\ge d-1$,
repeat the argument on those $p$ coordinates. Each coefficient
$p-j+1$ that must be cancelled is positive, hence nonzero in $\C$.
Eventually all first derivatives of $e_{p,1}$ would vanish, an
impossibility. The converse follows from the degree of each cofactor.
Thus the critical codimension is $N-d+2$, and
Corollary~\ref{cor:codim} proves the lower bound for $d\le N-1$.
For $d=N$, degree gives $\bdc(e_{N,N})\ge N$ and the diagonal
matrix gives equality.

For the upper bound, use vertices $(j,k)$ with
$0\le j\le d$ and $0\le k\le N-d$. The edge from $(j,k)$ to
$(j+1,k)$ has weight $x_{j+k+1}$, and the edge to $(j,k+1)$ has
weight $1$, whenever the indicated target exists. Source-to-sink
paths correspond bijectively to $d$-subsets of $[N]$. Their weight
sum is therefore $e_{N,d}$. Let $A$ be the weighted adjacency matrix
in topological order, with rows indexing sources and columns indexing
targets. It is strictly upper triangular. If $s,t$ denote the source
and sink, the affine-linear matrix
\[
 \begin{pmatrix}I-A&e_t\\-e_s^{\mathsf T}&0\end{pmatrix}
\]
has determinant $e_s^{\mathsf T}(I-A)^{-1}e_t=e_{N,d}$, by the Schur
complement and $\det(I-A)=1$. Its size is
$(d+1)(N-d+1)+1$. These estimates give absolute constants in both
Theta statements, including $d=N$.
\end{proof}

\section{Scope and limitations}

The proof is symbolic; finite computation is not a substitute for
any of its asymptotic steps. An accompanying exact checker verifies
selector minors, complete specialized permanent identities, small
instances of the partition identity behind~\eqref{eq:critical},
multihomogeneous coefficient expansions, and the integer parameter
estimates. It also verifies the Schur-complement kernel identities.
For controls, a smooth four-variable quadric represented by $\det_2$
has polar count $2=B(2,4)$; a smooth four-variable cubic section of
$\det_3$ has count $12\le B(3,4)=16$.

The full $n\times n$ determinant has critical codimension four, so
Corollary~\ref{cor:codim} applied directly to it gives only a linear
bound in $n$, consistent with its size-$n$ representation. In
particular, the smoothness hypothesis must not be omitted.

The result concerns determinantal representations, a model closely
related to algebraic branching programs. It supplies no lower bound
of the same order for unrestricted arithmetic circuits. For example,
$\sum_{i=1}^N x_i^N$ has a circuit of size $O(N\log N)$, yet
Proposition~\ref{prop:polar} gives border determinantal complexity
$\Omega(N^2)$. This directly precludes charging the same measure
linearly to arbitrary circuit gates.

The obstruction persists for syntactically multilinear circuits, even
without repeated squaring. For $N$ a power of two and $d=N/2$,
the elementary-symmetric argument supplies a smooth-slice product
$(N-d+1)(d-1)=\Theta(N^2)$. Yet $e_{N,d}$ has a division-free
syntactically multilinear circuit of size $O(N\log^2 N)$. Recursively
multiply the polynomials $1+x_i z$ by fast Fourier convolution on
disjoint variable blocks. A block of $b$ variables costs $O(b\log b)$
gates; summing over the balanced recursion proves the bound. Every
nonconstant multiplication has disjoint input supports. This standard
construction is recalled in~\cite[Section 5]{Shastri}. It rules out
even a charge weakened by any fixed polylogarithmic factor. Its
$O(N\log N)$ non-skew products lie outside the $O(N)$ budget of
Appendix~\ref{app:nonskew}; its Fourier gates mix degrees, so this
example alone makes no assertion about homogeneous internal gates.

Paired variables remove that remaining distinction for the present
\emph{affine} slice measure. Define
\[
 H_{N,d}(x,y)=[z^d]\prod_{i=1}^N(y_i+x_i z).
\]
The same FFT construction is now homogeneous and syntactically
set-multilinear in the groups $\{x_i,y_i\}$. In a recursion block $I$,
every coefficient and every Fourier value has degree $|I|$ and uses
exactly the groups in $I$. Additions preserve that group set, while
products combine disjoint child blocks. Zero padding may be removed
by constant propagation. The size remains $O(N\log^2 N)$. Setting
all $y_i=1$ gives $e_{N,d}$, so at $d=N/2$ the same quadratic
smooth-slice product is obtained by an affine restriction. Thus
homogeneity and syntactic set-multilinearity together still do not
repair a linear or polylogarithmically weakened charge of this
affine-slice measure to circuit size. A measure using only homogeneous
linear restrictions would be different and would not automatically
admit the degree-lowering permanent specialization used above.

There is also an intrinsic limit to this particular slice method.
An affine restriction of $\per_n$ has degree at most $n$. A smooth
homogeneous restriction of degree at least two can have at most $n^2$
variables: an extra kernel direction of the affine map would give a
projective critical point. Thus the product of degree and smooth-slice
dimension is at most cubic in $n$. The method cannot by itself prove
superpolynomial determinantal complexity, let alone VP$\ne$VNP.
Sheshadri~\cite{SheshadriReach} studies the broader degree-times-variables
scale of conormal intersection bounds. Here the number of variables of
the smooth target is $k=\Theta(n^2)$, while its degree is
$d=\Theta(\log n)$. Thus the present $\Theta(kd)$ scale is compatible
with that discussion; the quadratic scale for forms whose degree equals
their number of variables is a different specialization.

The supplementary verification code checks finite instances of the
identities and parameter inequalities. The uniform statements follow
from the proofs above and in the appendices.

\paragraph{Acknowledgment.}
The author thanks OpenAI's GPT-6 for assistance with exploratory
calculations, proof development, code checks, and manuscript drafting.
The author takes responsibility for the mathematical claims and the
final manuscript.

\appendix
\section{Matching identities and coefficient compression}\label{app:matching}

The finite-module argument and selector reproduce the ingredients
of~\cite[Chapter 5]{OAI} used above. Between them we give the
coefficient refinement, including its relation to~\cite{Guo}.

\subsection{The critical-locus dimension bound}

Put $q=d-1<t$. For each nonempty $B\subseteq[t]$, introduce a shared
parameter $p_B$. For a column vector $u\in\C^t$, define
\begin{equation}\label{eq:partitionQ}
 Q_i(u;p)=\sum_{\substack{I\subseteq[t]\setminus\{i\}\\|I|=q}}
       \ \sum_{\pi\vdash I}\mu(\pi)
                  \prod_{B\in\pi}(p_B-u_B),
 \quad
 u_B=\prod_{j\in B}u_j,
\end{equation}
where
$\mu(\pi)=\prod_{B\in\pi}(-1)^{|B|-1}(|B|-1)!$.
At a matrix $X$, specialize
$p_B=\sum_{a=1}^s\prod_{j\in B}x_{ja}$.
Inclusion--exclusion over equalities among column choices gives
\[
 Q_i(X_{\bullet a};p(X))=\partial_{ia}M_{t,s,d}(X).
\]
Indeed, $p_B-u_B$ sums the products belonging to columns other than
$a$, and the partition-lattice Mobius coefficient enforces that the
$q$ selected rows use distinct such columns.

When the $p_B$ are treated as coefficients, the highest $u$-degree
part of $Q_i$ is
\begin{equation}\label{eq:leadingQ}
 \alpha\,e_q(u_1,\ldots,\widehat{u_i},\ldots,u_t),
 \qquad \alpha=(-1)^q q!.
\end{equation}
To see the coefficient, choose $-u_B$ in every factor of
\eqref{eq:partitionQ}. For each fixed $I$, the sum of
$\prod_{B\in\pi}(|B|-1)!$ over partitions is $q!$: it counts
permutations by their cycle partitions. Every other term has
$u$-degree less than $q$.

Take $s$ column variables $u^{(a)}$ and, for each column, base
parameters $c_{j,a}$ for $1\le j<q$. Work over the polynomial ring
\[
 A=\C[p_B,\ c_{j,a}:\ \varnothing\ne B\subseteq[t],\
                         1\le j<q,\ 1\le a\le s].
\]
Form the quotient of $A[u^{(1)},\ldots,u^{(s)}]$ by
$Q_i(u^{(a)};p)=0$ and $e_j(u^{(a)})=c_{j,a}$ for $j<q$.
We claim that this quotient is a finite $A$-module.

Fix a column. Summing~\eqref{eq:leadingQ} over $i$ and using
$\sum_i e_q(u_{\widehat i})=(t-q)e_q(u)$ expresses $e_q(u)$
as a polynomial of $u$-degree below $q$, with coefficients in $A$.
Here $\alpha(t-q)\ne0$. The identity
\[
 e_q(u_{\widehat i})=
       \sum_{j=0}^q(-u_i)^j e_{q-j}(u)
\]
then turns each equation $Q_i=0$ into a monic relation
$u_i^q=$ terms of total $u$-degree below $q$; the intermediate
$e_{q-j}$ have been replaced by their base parameters.
Reducing these powers strictly lowers total column degree.
Doing this for every column proves that the finitely many monomials
with all exponents less than $q$ generate the quotient over $A$.

Its Krull dimension is consequently at most
$\dim A=2^t-1+s(q-1)$. The coordinate ring of the critical locus
is a quotient: substitute $u^{(a)}=X_{\bullet a}$,
$p_B=p_B(X)$, and $c_{j,a}=e_j(X_{\bullet a})$. All defining
relations hold, and the map is surjective because it includes every
$x_{ia}$. Dimension cannot increase under this quotient. This proves
\eqref{eq:critical}. In particular, the shared parameters need not
be independent at a critical point; enlarging to independent base
parameters only enlarges the ambient incidence.

\subsection{The coefficient image and its constant block}\label{app:compression}

We refine the preceding argument by counting the coefficients of
$Q_i$, rather than all the moments. The same disjointness-matrix
coefficient count is used in~\cite[Section 3]{Guo}. The additional
identity~\eqref{eq:constant-block} removes its constant block on
the critical locus. We give the details over an algebraically closed
field of characteristic zero or characteristic $p>\max(t,s)$.

For $\varnothing\ne S\subseteq[t]$ with $|S|\le d-1$, put
\[
 h_S=\sum_{\pi\vdash S}\mu(\pi)\prod_{B\in\pi}p_B,
 \qquad h_\varnothing=1.
\]
For $\varnothing\ne T\subseteq[t]$, $|T|\le d$, define
\[
 c_T=\sum_{\substack{S\subseteq[t]\setminus T\\|S|=d-|T|}}h_S.
\]
Collecting coefficients in~\eqref{eq:partitionQ} gives
\begin{equation}\label{eq:collected-Q}
 Q_i(u;p)=\sum_{\substack{J\subseteq[t]\setminus\{i\}\\|J|\le d-1}}
       (-1)^{|J|}|J|!\,c_{J\cup\{i\}}u_J.
\end{equation}
To see this, choose the blocks supplying a factor $-u_B$ and let
$J$ be their union. The partitions of the remaining rows contribute
$h_S$. On $J$, summing the products of cyclic-order counts over its
partitions gives $|J|!$, with sign $(-1)^{|J|}$. Summing over the
remaining $S$ proves the identity. The leading coefficients, for
$|J|=d-1$, have $c_{J\cup\{i\}}=1$.

For $|T|=j$, the vector $(c_T)$ is a linear image of $(h_S)_{|S|=d-j}$
under the disjointness matrix on $j$-subsets and $(d-j)$-subsets
of $[t]$. Its rank is
\[
 \min\left\{\binom tj,\binom t{d-j}\right\}
 =\binom t{\min(j,d-j)}.
\]
Here is a proof of the rank fact, including its characteristic range.
Transpose and complement subsets to reduce to the following assertion:
if $a\le v\le t-a$ and
\[
 \sum_{\substack{A\subseteq V\\|A|=a}}z_A=0
 \quad\hbox{for every $v$-subset $V\subseteq[t]$},
\]
then every $z_A=0$. Induct on $a$, the case $a=0$ being immediate.
For distinct $i,j$, subtract the equations for $U\cup\{i\}$ and
$U\cup\{j\}$, where $U$ ranges over the $(v-1)$-subsets avoiding
$i,j$. On the remaining $t-2$ elements the induction hypothesis
applies to the coefficients $z_{C\cup\{i\}}-z_{C\cup\{j\}}$:
$a-1\le v-1\le(t-2)-(a-1)$. They all vanish. The exchange graph
of $a$-subsets is connected, so every $z_A$ has one common value $z$.
The original equation now reads $\binom va z=0$; this binomial
coefficient is nonzero in the stated characteristics. This proves
the rank formula, a form of the classical inclusion-matrix rank
theorem also used in~\cite{Guo}.

The change from the independent moments $p_S$ to the $h_S$ is
triangular by $|S|$, with nonzero diagonal coefficient
$(-1)^{|S|-1}(|S|-1)!$. Thus the $h_S$ are independent coordinates.
Different $j$ use different degrees $d-j$, and their images occupy
different monomials in~\eqref{eq:collected-Q}. Consequently the
formal coefficient image has dimension exactly
\[
 R(t,d)=\sum_{j=1}^{d-1}\binom t{\min(j,d-j)}.
\]
The $j=1$ block consists of $t$ independent constant coefficients.

At an actual matrix, inclusion--exclusion identifies $h_S$ with the
sum of the matchings using exactly the row set $S$. For each $i$,
direct counting gives
\begin{equation}\label{eq:constant-block}
 \sum_{a=1}^s\partial_{ia}M_{t,s,d}=(s-d+1)c_{\{i\}}.
\end{equation}
A matching of size $d-1$ avoiding row $i$ is counted once for each
unused column. Since $s-d+1\ne0$ in the ground field, every
$c_{\{i\}}$ vanishes on the critical locus. Imposing this zero
block on the formal coefficient image leaves dimension
$R_{\rm crit}(t,d)=R(t,d)-t$.

Choose coordinates on each remaining disjointness image and use
them to reconstruct the coefficients of~\eqref{eq:collected-Q}.
Together with the $s(d-2)$ column elementary invariants used above,
these coordinates generate a polynomial parameter ring of dimension
$R_{\rm crit}(t,d)+s(d-2)$. Its universal column incidence still
contains every critical matrix with its actual coefficient data.
The leading part of $Q_i$ is unchanged, so the same degree-reducing
monic relations make its coordinate ring finite over this parameter
ring. Surjecting onto the critical-locus coordinate ring proves
\eqref{eq:critical-compressed}.

For use in the parameter estimates, the count has the explicit forms
\[
 \begin{split}
 R_{\rm crit}(t,2h)&=2\sum_{j=1}^{h-1}\binom tj+\binom th-t,\\
 R_{\rm crit}(t,2h+1)&=2\sum_{j=1}^{h}\binom tj-t.
 \end{split}
\]
The dimension calculation is exact for the formal coefficient image.
It does not assert that actual critical matrices fill that image,
or that the resulting upper bound on their dimension is sharp.


\subsection{The selector identity}

Use the square-zero algebra
$\C[z_1,\ldots,z_r]/(z_1^2,\ldots,z_r^2)$ and set
$y_h=\sum_{i\in R_h}z_i$. Then $y_h^{t+1}=0$. The product of
the linear forms associated with the constant columns of $U$ is
\begin{equation}\label{eq:selectorproduct}
 \mathcal S:=\prod_{h=1}^b y_h^{t-d}
 \prod_{h=2}^b
       \bigl((y_1+\cdots+y_{h-1})^d+\theta y_h^d\bigr),
 \qquad \theta=(-1)^{d+1}2^d.
\end{equation}
This follows from
$\prod_{j=0}^{d-1}(S+2\zeta^jY)=S^d+\theta Y^d$.

In this truncated algebra, induction on the number of blocks gives
\begin{equation}\label{eq:selectorcollapsed}
 \mathcal S
   =\sum_{h=1}^b\lambda_h y_h^{t-d}\prod_{j\ne h}y_j^t.
\end{equation}
For the induction, every term has exactly one block of exponent
$t-d$, with all other old blocks of exponent $t$. In the next
factor, $(y_1+\cdots+y_{h-1})^d$ can survive only by taking all
$d$ powers from that one unsaturated block. It saturates the old
block and leaves the new one unsaturated. The other term,
$\theta y_h^d$, keeps the old unsaturated block and multiplies
its coefficient by $\theta$. Thus old coefficients are multiplied
by $\theta$, and the new coefficient is their previous sum.
Starting from $1$, their sums are $(1+\theta)^{h-1}$, yielding
the displayed formulas for $\lambda_h$ in the main text.

For a set $I$ of $d$ omitted rows, the coefficient of
$\prod_{i\notin I}z_i$ in the product of the column forms is
$\per(U_{[r]\setminus I,\bullet})$. Equation~\eqref{eq:selectorcollapsed}
makes this zero unless $I$ lies within one block $R_h$. In that case
the coefficient is
$\lambda_h(t-d)!(t!)^{b-1}$, independent of the particular $I$.
Finally expand the full permanent along its $s-d$ all-ones bottom
rows. Exactly $d$ columns from the $X$ block remain for the top rows,
and the bottom contribution is $(s-d)!$. Summing the remaining
size-$d$ matchings gives
\[
 \per\begin{pmatrix}X&U\\\mathbf1&0\end{pmatrix}
  =(s-d)!(t-d)!(t!)^{b-1}
      \sum_{h=1}^b\lambda_h M_{t,s,d}(X^{(h)}),
\]
as required. Both $\theta$ and $1+\theta$ are nonzero, so every
weight used in the critical-locus product is nonzero.

\section{Algebraic transfer with polar multiplicities}\label{app:algebraic}

This appendix gives a second proof of the smooth-target transfer and
records its characteristic restrictions. The determinantal incidence
and its multidegree are still the credited mechanism of
\cite{Sheshadri}. The additional argument below replaces persistence
of reduced complex points by conservation of isolated
complete-intersection length. In particular, it does not assume that
the Gauss map is separable.

\begin{lemma}[Local conservation]\label{lem:conservation}
Let $R$ be a complete discrete valuation ring with algebraically
closed residue field $K$ and uniformizer $t$. Let
$F_1,\ldots,F_k\in R[x_1,\ldots,x_k]$. Suppose their special fiber
has a finite set of isolated points $z$, with local lengths $\mu_z$.
There are finite flat neighborhoods of these points in the common
zero scheme, of combined rank $\sum_z\mu_z$ over $R$. Any function
nonzero at a selected special point is a unit on its corresponding
finite local factor.
\end{lemma}

\begin{proof}
Take the quasi-finite open locus around the selected points.
Zariski's Main Theorem embeds it into a finite $R$-scheme. Since
$R$ is henselian, this finite scheme splits into factors indexed by
its closed-fiber points. A factor whose closed point lies in the
chosen open lies entirely in that open: every nonempty closed
subset of a finite local factor contains its closed point. These
factors give finite neighborhoods in the original zero scheme.

At each selected point the ambient local ring is regular of
dimension $k+1$. The $k$ equations followed by $t$ form a system
of parameters, because their quotient has finite length. A system
of parameters in a regular local ring is a regular sequence.
Thus $t$ is a non-zero-divisor after imposing the $F_i$. Each
finite local factor is torsion-free, hence free, over $R$, with
rank equal to its special-fiber length. The unit assertion follows
by reduction to the closed point. These are standard
complete-intersection and henselian arguments; compare
\cite[Sections 10.104, 10.153 and Lemma 93.17.4]{Stacks}.
\end{proof}

\begin{proposition}\label{prop:anychar}
Let $K$ be algebraically closed, in any characteristic. If $f$ is
a homogeneous degree-$d$ form in $k\ge3$ variables and
$V(f)\subset\PP^{k-1}_K$ is smooth, then
\[
 \bdc_K(f)\le m\quad\Longrightarrow\quad
 d(d-1)^{k-2}\le B(m,k).
\]
Consequently $\bdc_K(f)\ge(k-1)(d-1)/(4e)$.
\end{proposition}

\begin{proof}
Set $X=V(f)$. Smoothness implies that the partial derivatives
have no common zero on $X$. They define a morphism
\[
 \gamma:X\longrightarrow\PP^{k-1},\qquad
 \gamma^*\mathcal O(1)=\mathcal O_X(d-1).
\]
This morphism is finite onto its image: a positive-dimensional
fiber would contain a projective curve on which the ample bundle
$\mathcal O_X(d-1)$ had degree zero. Hence $k-2$ general
hyperplanes in the target pull back to a finite scheme of length
\[
 \delta=\int_X c_1(\mathcal O_X(d-1))^{k-2}
        =d(d-1)^{k-2}.
\]
It is cut out by directional derivatives $D_1f,\ldots,D_{k-2}f$.
Choose $\ell$ nonzero on its finite support. In the affine
normalization $\ell=1$, the equations are
\[
 f=0,\quad \ell-1=0,\quad D_1f=\cdots=D_{k-2}f=0.
\]
This is a zero-dimensional complete intersection of length
$\delta$. It can be nonreduced. At every point, however, some
first partial of $f$ is nonzero. We have used the gradient only
on $X$; no division by $d$ in Euler's identity is needed.

The affine size-$m$ determinants form a constructible image in
coefficient space. Algebraic curve selection at a point of its
closure, followed by a finite function-field extension, normalization,
and completion, gives a complete discrete valuation ring $R$ with
residue field $K$ and a polynomial
\[
 g\in R[x_1,\ldots,x_k],\qquad g\bmod t=f,
 \qquad g=\det A\text{ over }\operatorname{Frac}(R).
\]
The matrix entries may have poles; only the coefficients of $g$
are required to be regular. A constant family suffices when $f$
already has an exact representation.

Apply Lemma~\ref{lem:conservation} to $g,\ell-1,D_1g,\ldots,D_{k-2}g$.
It retains a generic-fiber scheme of length $\delta$. A first
partial nonzero at each selected special point stays a unit on
its finite factor. Thus every retained generic matrix value has
corank one. Extend the generic field to an algebraic closure,
which does not change length, and choose $\Lambda$ injective on
all of the finitely many matrix images.

The Schur-complement elimination from Section~3 is
scheme-theoretic: it divides only by units and identifies the
kernel incidence locally with $g=0$. The identity
$\adj(A)=\det(B)vu^{\mathsf T}$ identifies the lifted polar
equations with $D_jg/\det(B)$. Hence it preserves the complete
local polar schemes, including their nilpotents and total length.

Finally perturb the multihomogeneous incidence equations by a
new formal parameter to general sections of the same globally
generated line bundles. Their proper intersection has length
\[
 [H^k U^{m-1}V^{m-1}]
 H(H+U)^m(H+V)^{m-1}(U+V)^{k-2}=B(m,k).
\]
Apply Lemma~\ref{lem:conservation} locally on affine charts of the
smooth product of projective spaces. Every original isolated
component retains its length. Thus $\delta\le B(m,k)$ even
when the original incidence has excess components elsewhere.
The numerical estimate in Section~3 is characteristic-independent.
\end{proof}

\begin{corollary}[Large characteristic]\label{cor:largechar}
Let $K$ be algebraically closed of characteristic $p>n$.
Both inequalities of Theorem~\ref{thm:main} hold over $K$ in their
respective ranges $n\ge64$ and $n\ge512$.
\end{corollary}
\begin{proof}
The matching finite-module proof divides by $(d-1)!$ and $t-d+1$.
The compression additionally uses $s-d+1$ and the ranks of inclusion
matrices on a $t$-element set. All these arguments apply for
$p>\max(t,s)$, by Appendix~\ref{app:compression}.
The selector uses roots of unity of order $d<p$ and the nonzero factor
$(s-d)!(t-d)!(t!)^{b-1}$. For both parameter choices in the proof of
Theorem~\ref{thm:main}, $t,s\le n$ and $2^d<n$; hence $\theta$ and
$1+\theta$ are nonzero modulo $p$. The smooth-slice incidence argument
applies because $p>d$, so Euler's identity identifies the critical
cone with the hypersurface singular cone. Proposition~\ref{prop:anychar}
and the same integer estimates prove the claim.
\end{proof}

This is not a superquadratic result in any fixed positive
characteristic. In particular, removing the analytic transfer
does not remove the selector's factorial obstruction.

\paragraph{Multiplicity controls.}
For $d=p+1$, the smooth Fermat form in $k$ variables has
inseparable polar equations $x_i^p=0$ on a suitable affine chart.
The polar length is $(p+1)p^{k-2}$, whereas its support has only
$p+1$ points. The triangular deformation $x_i^p+t=0$ retains
this length. A smooth characteristic-two conic also checks the
case $p\mid d$, where the ambient gradient can have a base point
outside the hypersurface. These exact finite controls test why
length, rather than distinct-point count, is needed; the proof
above supplies the general assertion.

\section{A charge for non-skew multiplication gates}\label{app:nonskew}

This appendix extends the polar count to a circuit model with a stated
budget on non-skew multiplication gates. A binary multiplication
gate is called \emph{non-skew} if neither child is an input variable or
field constant. Other multiplication gates are skew. Circuits are
commutative and division-free; arbitrary fan-out and cancellation are
allowed. The related noncommutative restrictions studied
in~\cite{LMS} are different models, and their lower bounds are not used.

\begin{theorem}[Circuit-state polar count]\label{thm:nonskew}
Let $f$ be a homogeneous degree-$d$ form in $k\ge3$ variables with
smooth projective hypersurface. Suppose $f$ is a fixed affine
restriction of a coefficient limit of outputs of circuits with $S$
internal binary gates, $N$ input variables, and at most $q$ non-skew
multiplication gates, where $0\le q\le S$. A larger stated budget may
always be replaced by $S$. Put $M=S+N+1$. Then
\[
 d(d-1)^{k-2}\le T_q(M,k),
\]
where $T_0(M,k)=B(M,k)$ and, for $q\ge1$,
\begin{equation}\label{eq:nonskew-count}
 T_q(M,k)\le
 2^{q+k-2}\binom{q+k-3}{k-2}\binom{2M-1}{k-1}.
\end{equation}
In particular, for every fixed $C\ge0$, the budget $q\le C(k-1)$
implies $M=\Omega_C(kd)$ as $d$ tends to infinity.
\end{theorem}

\begin{corollary}\label{cor:nonskew-per}
For each fixed constant $C$, a division-free circuit computing
$\per_n$ with at most $Cn^2$ non-skew multiplication gates has
$\Omega_C(n^2\log n)$ internal gates. The same holds for coefficient
approximation with these size and non-skew-gate budgets.
\end{corollary}

\begin{proof}
Use the fixed smooth restriction from the proof of
the first part of Theorem~\ref{thm:main}, taking
$t=\lfloor\log_2 n\rfloor$, $d=\lfloor t/2\rfloor$,
$k=\lfloor n^2/16\rfloor$ for all sufficiently large $n$,
and $N=n^2$ original
input registers. Affine restriction changes their labels to affine
forms but preserves their role as skew operands in the construction
below. Theorem~\ref{thm:nonskew} gives
$S+N+1=\Omega_C(n^2\log n)$; subtracting $N+1$ proves the assertion
for all sufficiently large $n$.
\end{proof}

\subsection*{State equations and their adjoints}

Introduce one state $w$ for each original input variable and each
internal gate, for a total of $M-1$ states. After restriction, the input
equations are $w_i=L_i(x)$ with $L_i$ affine. In topological order, each
addition, subtraction, or skew multiplication has an equation affine
in the states, with coefficients affine in $x$. For a skew product,
replace its input child by the affine label $L_i(x)$; retain the state
register for the other operand. Constants require no state register.
Each non-skew product has an equation
\[
 w_g-w_aw_b=0.
\]
The Jacobian of these $M-1$ graph equations with respect to the states
is triangular with diagonal one. Thus they define the graph of the
computation, including every use of every shared value.

Add the output equation $w_{\rm out}=0$. On
$\PP^k_{[t:x]}\times\PP^{M-1}_{[v_0:w]}$, the $M-q$ linear-state
equations, including the output, become forms $F_i$ of bidegree
$(1,1)$. Examples are
\[
 tv_i-\widehat L_i(x)v_0,\qquad
 tv_g-\widehat L_i(x)v_a,\qquad tv_{\rm out}.
\]
For the other $q$ equations use
\[
 tQ_j(v)=t(v_gv_0-v_av_b),
\]
of bidegree $(1,2)$. The added factors $t$ have no effect on the
chart $t=v_0=1$. Their extra components at infinity are permitted.

At a zero of the output with nonzero differential, the Jacobian of
all $M$ constraints in the $M-1$ state directions has rank $M-1$.
Its left kernel is a unique multiplier line. Write the multipliers as
$u_i$ for the $F_i$ and $\lambda_j$ for the $tQ_j$. The multiplier of
the output equation is nonzero: otherwise the invertible graph
Jacobian forces every multiplier to vanish. Normalize that multiplier
to one. Implicit differentiation of the graph equations shows that
\[
 S_y=\sum_i u_i(F_i)_y+\sum_j\lambda_j(tQ_j)_y
\]
vanishes for all state directions, and its $x$-components are, up to
sign, the differential of the output. The state adjoints sum all
contributions from fan-out; no parse-tree expansion occurs.

Choose a simple polar fiber as in Proposition~\ref{prop:polar}. It
has $\delta=d(d-1)^{k-2}$ points. At each point the graph equations
eliminate the states and the adjoint equations eliminate the normalized
multipliers, by invertible Jacobians. The resulting local system is
the original square polar system. Every lifted point is therefore
isolated and reduced. The same is true for the $\delta$ points
persisting in any sufficiently close polynomial circuit output.

\subsection*{A bundle count that respects reuse}

Write $H,V$ for the hyperplane classes of
$Y=\PP^k\times\PP^{M-1}$ and use the quotient projective bundle
\[
 \pi:\PP(E')\longrightarrow Y,\qquad
 E'=\mathcal O_Y^{\oplus(M-q)}\oplus
       \mathcal O_Y(V)^{\oplus q},\qquad
 L=c_1(\mathcal O_{\PP(E')}(1)).
\]
The coordinates $u_i$ have class $L$, whereas the $\lambda_j$ have
class $L-V$. Hence a state-adjoint equation has class $L+H$, and an
$x$-directional adjoint equation has class $L+V$. The latter has no
quadratic-constraint contribution, since $Q_j$ is independent of $x$;
it is nevertheless a section of the stated line bundle. The
homogenizing coordinate $t$ is not one of these $x$ directions.

The incidence has $M-q$ equations of class $H+V$, $q$ of class
$H+2V$, $M-1$ state adjoints of class $L+H$, $k-2$ polar adjoints of
class $L+V$, and one equation $\ell(x)-t$ of class $H$. Their number
is $k+2M-2=\dim\PP(E')$. All these bundles are globally generated:
$E'$ is globally generated, as are its twists by $H$ and $V$.
General perturbations of the sections thus have proper
zero-dimensional intersection. The isolated points just constructed
persist, giving the upper bound
\begin{equation}\label{eq:state-intersection}
 T_q(M,k)=\int_{\PP(E')}
 H(H+V)^{M-q}(H+2V)^q(L+H)^{M-1}(L+V)^{k-2}.
\end{equation}
This bound concerns the original system's isolated points and allows
excess components elsewhere. In particular, it does not assume
properness of the incidence at infinity.

For $q=0$ the bundle is trivial and
\eqref{eq:state-intersection} is $B(M,k)$ after exchanging the two
kernel-factor names. For $q\ge1$, the projective-bundle formula gives
\[
 \pi_*(L^{M-1+a})=\binom{q+a-1}{a}V^a\quad(a\ge0),\qquad
 \pi_*(L^j)=0\quad(j<M-1).
\]
For precision about conventions, these formulas follow from
$L^{M-q}(L-V)^q=0$ and $\pi_*(L^{M-1})=1$ by induction;
see~\cite[Sections 27.21 and 42.36]{Stacks}. No sign convention for
Segre classes is implicit.

All coefficients in the integrand and pushforward are nonnegative.
Bound $(H+2V)^q$ coefficientwise by $2^q(H+V)^q$. If $i$ copies
of $H$ are chosen from $(H+V)^M$ and $c$ from $(L+H)^{M-1}$, the
required $H$ degree is equivalent to $i+c=k-1$. Choosing $j$ copies
of $L$ from $(L+V)^{k-2}$ leaves a pushforward contribution bounded by
\[
 \sum_{j=c}^{k-2}\binom{k-2}{j}\binom{q+j-c-1}{j-c}
 \le 2^{k-2}\binom{q+k-3}{k-2}.
\]
The remaining sum is
\[
 \sum_{i+c=k-1}\binom Mi\binom{M-1}c
 =\binom{2M-1}{k-1},
\]
which proves~\eqref{eq:nonskew-count}. The count is nondecreasing
in $q$: both the constraint coefficients and the pushforward
coefficients increase. An upper budget on $q$ therefore suffices.

Fixed-size circuit families form a finite union of constructible
polynomial images. The simple target polar points persist under
coefficient approximation, with no boundedness assumption on circuit
parameters or state values. This proves the border assertion and
completes the proof of Theorem~\ref{thm:nonskew}.

\subsection*{An explicit tradeoff and its limit}

Taking $(k-1)$-st roots and using the standard binomial estimates
gives the convenient bound, for all $q\ge0$,
\begin{equation}\label{eq:state-tradeoff}
 M\ge\frac12+
 \frac{(k-1)(d-1)}
 {4e^2\,2^{q/(k-1)}(1+q/(k-2))}.
\end{equation}
For $q=0$, Proposition~\ref{prop:polar} gives a stronger constant.
For $q\ge1$, use
\[
 \binom{q+k-3}{k-2}^{1/(k-1)}
 \le e\left(1+\frac q{k-2}\right),\qquad
 \binom{2M-1}{k-1}^{1/(k-1)}
 \le\frac{e(2M-1)}{k-1}.
\]
If $q\le C(k-1)$, the denominator of
\eqref{eq:state-tradeoff} is bounded by a constant depending only on
$C$, proving the last assertion of the theorem.

This charge tolerates $O(k)$ non-skew gates but weakens as their number
grows. For example, repeated squaring computes
$\sum_{i=1}^k x_i^{2^a}$ using $k(a-1)$ non-skew squarings after the
first squares. Here $q/k$ grows with $\log d$, outside the fixed-$C$
regime. Thus the argument respects the earlier counterexample to a
linear charge of smooth-slice score to unrestricted circuit size.
Corollary~\ref{cor:nonskew-per} is a restricted circuit result, not an
unrestricted $\Omega(n^2\log n)$ circuit bound.

The accompanying finite checker independently expands 416 intersection
coefficients, checks 364 monic projective-bundle reductions and 32
repeated-squaring controls, and verifies 35 full local incidence
Jacobians, including 77 multiply-used internal states. These checks
support the displayed identities and local elimination; the preceding
argument proves the uniform statement.

\section{Generic output levels and valid divisions}\label{app:rationalstate}

This appendix extends the state construction to valid divisions and
formal rational-output approximation over $\C$. It is separate from the
main determinantal theorem and retains an explicit gate budget. In
particular, it does not extend the two-kernel determinant incidence to
nonzero determinant levels.

\begin{theorem}\label{thm:rationalstate}
Let a circuit over $\C((\varepsilon))$ have $S$ internal binary gates,
$N$ inputs, arbitrary fan-out, and valid divisions. Replace divisions
by nonzero field constants by scalar multiplications. Count in $q$
every remaining division and every non-skew multiplication, as defined
in Appendix~\ref{app:nonskew}; thus $0\le q\le S$. Suppose its rational
output satisfies
\[
 r_\varepsilon(X)=F(X)+O(\varepsilon)
 \quad\hbox{in }\C(X)((\varepsilon)),
\]
where $F$ is a polynomial. Suppose a fixed affine restriction $F(Ax)$
is a homogeneous degree-$d$ form $f$ in $k\ge3$ variables, with $d\ge2$
and smooth projective hypersurface. Set $M=S+N+1$. Then
\begin{equation}\label{eq:rationalstate}
 d(d-1)^{k-1}\le T_q(M,k+1).
\end{equation}
Consequently,
\begin{equation}\label{eq:rationaltradeoff}
 M\ge\frac12+
 \frac{k(d-1)}{4e^2\,2^{q/k}(1+q/(k-1))}.
\end{equation}
For every fixed $C$, the budget $q\le Cn^2$ therefore forces
$S=\Omega_C(n^2\log n)$ for such an approximation to $\per_n$.
\end{theorem}

Here $A$ denotes an affine map, including its constant term. The
output need not be polynomial for nonzero parameter. The gate graph
and its input roles are fixed; constants may be Laurent series.

\subsection*{A full-dimensional polar target}

Smoothness and Euler's identity imply $\Crit(f)=\{0\}$. The gradient
defines a finite morphism
\[
 \gamma:\PP^{k-1}\longrightarrow\PP^{k-1},\qquad
 \gamma^*\mathcal O(1)=\mathcal O(d-1),
\]
of degree $(d-1)^{k-1}$. In characteristic zero it is generically
\'{e}tale. A general direction $a$ avoids both its branch locus and
$\gamma(V(f))$. Its projective fiber has $(d-1)^{k-1}$ distinct
points at which $f$ is nonzero. For each nonzero $c$, every such point
has exactly $d$ scalar representatives with $f=c$. Thus the map
\[
 \Psi_f:x\longmapsto(f(x),[\nabla f(x)])
\]
has a general fiber of $\delta=d(d-1)^{k-1}$ simple points. In a target
chart $a_k=1$ the corresponding square system is
\begin{equation}\label{eq:levelpolar}
 f-c=0,\qquad f_{x_i}-a_i f_{x_k}=0\quad(1\le i<k).
\end{equation}
Its Jacobian is invertible: the direction equations are \'{e}tale in
projective coordinates, while the remaining radial derivative is
$d\,c$, which is nonzero.

For any nonzero polynomial $E_0$, the closure of the image of
$\{E_0=0\}$ under this map, on its domain, has dimension at most
$k-1$. The target has dimension $k$. A general $(c,a)$ consequently
has all its $\delta$ preimages outside $\{E_0=0\}$. Allowing the
output level to vary is essential: a circuit with output $f^2/f=f$
is invalid at every point of the zero level.

\subsection*{A ramified translation of the fixed slice}

We give the argument because direct substitution of a fixed slice
can annihilate internal divisors. Represent each intermediate value
as a quotient of polynomials in $X$ over $\C((\varepsilon))$.
There is a finite list of nonzero polynomials whose nonvanishing
ensures every division is valid; include all these numerators and
denominators, and an output denominator. Normalize each polynomial
$H$ by a Laurent scalar so that $H\in\C[[\varepsilon]][X]$ and its
reduction $H_0$ is nonzero. The degrees are finite. Choose a constant
vector $b$ such that no $H_0(A(0)+\tau b)$ is identically zero, and
choose an integer $R$ larger than all these degrees. These choices
exist because there are finitely many proper exceptional algebraic
sets of vectors $b$.

Substitute $\varepsilon=\tau^R$ and $X=Ax+\tau b$. The $\tau$-adic
valuation of each $H_0(Ax+\tau b)$ is at most $\deg H_0<R$.
Terms of positive $\varepsilon$ order have $\tau$ order at least
$R$, so cannot cancel this leading term. Every original division
therefore remains valid as a rational function after substitution.
Neither $S$ nor $q$ changes: an original input operand is still
represented by its affine label in the state construction.

Write the output as $N/D$, normalizing the denominator to have
coefficientwise valuation zero. The approximation hypothesis and
the Gauss valuation imply $N-FD\in\varepsilon\C[[\varepsilon]][X]$.
If $a<R$ is the valuation of the substituted denominator, write it
as $\tau^a E$, with $E\in\C[[\tau]][x]$ and $E_0\ne0$. The new
rational output has the form
\begin{equation}\label{eq:ramifiedoutput}
 g(\tau,x)=F(Ax+\tau b)+\tau^{R-a}\frac{H(\tau,x)}{E(\tau,x)},
 \qquad H\in\C[[\tau]][x].
\end{equation}
On $E_0\ne0$, the output and its $x$-derivatives reduce to $f$ and
its derivatives.

Choose $(c_0,a_0)$ so that the $\delta$ simple points of
\eqref{eq:levelpolar} avoid $E_0=0$. For every $b'\in\C^k$,
the target $(c_0,a_0)+\tau b'$ has $\delta$ distinct simple Hensel
lifts for the system with $f$ replaced by $g$. At these lifts $E$
and $g_{x_k}$ are units. Internal canceled poles still require a
separate check. Over $K=\C((\tau))$, let $U$ be the nonempty open
set where all original divisions are valid and let
$V=\{E\ne0,\ g_{x_k}\ne0\}$. On $V$ the map
\[
 \Psi_g=(g,g_{x_1}/g_{x_k},\ldots,g_{x_{k-1}}/g_{x_k})
\]
is a morphism to $\mathbb A^k_K$. The closure of the image of
$V\setminus U$ has dimension at most $k-1$. The set
$(c_0,a_0)+\tau\C^k$ is Zariski dense over $K$: an invertible affine
change reduces this to the density of the infinite Cartesian set
$\C^k$ for polynomials over $K$. Choose $b'$ avoiding that bad image.
All the retained Hensel lifts then belong to $U$.

\subsection*{Shared states, division pivots, and the count}

Introduce the $M-1$ input and gate states as in
Appendix~\ref{app:nonskew}. A division $w_g=w_a/w_b$ has equation
$w_gw_b-w_a=0$. At every retained valid point the graph Jacobian in
topological state order is triangular, with pivot one at ordinary
gates and nonzero pivot $w_b$ at a division. It is invertible.
Add $w_{\rm out}-c(\tau)=0$. The state adjoints have a unique
multiplier line whose output multiplier is nonzero. Normalizing it
to one and eliminating states and adjoints recovers
\eqref{eq:levelpolar} with $g,c(\tau),a(\tau)$ in place of
$f,c,a$. Implicit differentiation includes every outgoing use of a
shared gate. The lifted incidence points are isolated and reduced.

On $Y=\PP^k_H\times\PP^{M-1}_V$, homogenize the $M-q$ linear-state
constraints to class $H+V$. The $q$ quadratic-state constraints have
class $H+2V$; for division use $t(v_gv_b-v_av_0)$. The added factor
$t$ is a unit on the chosen affine chart. Use the same quotient
bundle $E'=\mathcal O^{M-q}\oplus\mathcal O(V)^q$ and tautological
class $L$ as in Appendix~\ref{app:nonskew}. There are $M-1$ state
adjoints of class $L+H$ and now $k-1$ polar adjoints of class $L+V$.
There is no equation $\ell=1$. The total number of equations is
$k+2M-2=\dim\PP(E')$. Their bundles are globally generated, so the
same isolated-point conservation argument gives
\[
 \delta\le\int_{\PP(E')}
 (H+V)^{M-q}(H+2V)^q(L+H)^{M-1}(L+V)^{k-1}
 =T_q(M,k+1).
\]
For the last identity, the additional $H$ in
\eqref{eq:state-intersection} cancels the additional required base
$H$ degree when $k$ is replaced by $k+1$. All counts may be made
over an algebraic closure of $K$; no boundedness of gate constants
or state values is needed. Substituting $k+1$ into
\eqref{eq:nonskew-count} and the root estimate proves
\eqref{eq:rationaltradeoff}. Applying the fixed permanent slice and
subtracting $N+1$ proves the final assertion of the theorem.

The independent finite controls check 294 shifted coefficient
identities, 24 full rational-circuit incidence Jacobians, two
canceled-pole examples, and eight Fermat level counts. They reject
24 deliberately incorrect unit-pivot substitutions. These are
implementation checks, not a substitute for the argument above.
No extension of this generic-level proof to positive characteristic
is asserted here.


\begin{thebibliography}{99}\raggedright
\bibitem{LMR}
J. M. Landsberg, L. Manivel, and N. Ressayre,
\emph{Hypersurfaces with degenerate duals and the geometric complexity
theory program}, Comment. Math. Helv. 88, 469--484.
\href{https://doi.org/10.4171/CMH/292}{doi:10.4171/CMH/292}.
\href{https://arxiv.org/abs/1004.4802}{arXiv:1004.4802}.

\bibitem{OAI}
OpenAI, \emph{Ten Advances in Mathematics and Theoretical Computer
Science}, Chapter 5, Proposition 5.1,
Corollary 5.2, and Lemma 6.1.
\url{https://cdn.openai.com/pdf/ten-proofs-oai.pdf}.

\bibitem{Sheshadri}
K. Sheshadri, \emph{A near-quadratic lower bound on the border
determinantal complexity of $\sum_i x_i^n$ via conormal specialization}.
\href{https://arxiv.org/abs/2606.13628}{arXiv:2606.13628v1}.

\bibitem{SheshadriReach}
K. Sheshadri, \emph{The Exact Reach of Conormal Invariants in
Determinantal Complexity: a Quadratic No-Go Theorem}.
\href{https://arxiv.org/abs/2606.15970v1}{arXiv:2606.15970v1}.

\bibitem{Guo}
J. Guo, \emph{An $n^2\log\log n$ Lower Bound for Permanent Circuits
with Valid Division}.
\href{https://arxiv.org/abs/2609.29568v1}{arXiv:2609.29568v1}.

\bibitem{Orzel}
I. Orzel, \emph{Computing the Elementary Symmetric Polynomials in
Positive Characteristics}, version 3.
\href{https://arxiv.org/abs/2509.05009v3}{arXiv:2509.05009v3}.

\bibitem{CKSV}
P. Chatterjee, M. Kumar, A. She, and B. L. Volk,
\emph{Quadratic Lower Bounds for Algebraic Branching Programs and
Formulas}, Computational Complexity 31, article 8.
\href{https://doi.org/10.1007/s00037-022-00223-8}{doi:10.1007/s00037-022-00223-8}.
\href{https://arxiv.org/abs/1911.11793}{arXiv:1911.11793}.

\bibitem{Stacks}
The Stacks Project Authors, \emph{The Stacks Project}.
Sections \href{https://stacks.math.columbia.edu/tag/00N7}{10.104 (Cohen--Macaulay rings)},
\href{https://stacks.math.columbia.edu/tag/04GE}{10.153 (Henselian local rings)},
\href{https://stacks.math.columbia.edu/tag/01OA}{27.21 (Projective bundles)}, and
\href{https://stacks.math.columbia.edu/tag/02TV}{42.36 (Projective space bundle formula)};
Lemma \href{https://stacks.math.columbia.edu/tag/0E7W}{93.17.4 (finite flat smoothing)}.

\bibitem{LMS}
N. Limaye, G. Malod, and S. Srinivasan,
\emph{Lower Bounds for Non-Commutative Skew Circuits},
Theory of Computing 12, article 12, 1--38.
\url{https://theoryofcomputing.org/articles/v012a012/}.

\bibitem{Shastri}
P. Shastri, \emph{Lower Bounds for Noncommutative Circuits with Low
Syntactic Degree}, ITCS, LIPIcs 362, article 115.
\url{https://doi.org/10.4230/LIPIcs.ITCS.2026.115}.
\end{thebibliography}
\end{document}
